\clearpage
\VerdierPhysicalPageAnchor{0143}%
\begingroup
\small
\setlength{\parskip}{0pt}%
\setlength{\emergencystretch}{1.5em}%

\begingroup\itshape
\VerdierSemanticAnchor{ch2-2-3-1-continuation}%
Pour toute sous-catégorie triangulée pleine $\mathcal A$ telle que
$\mathcal B\subset\mathcal A\subset\mathcal D$, désignons par
$S_{\mathcal D}(\mathcal B)$ (resp. $S_{\mathcal A}(\mathcal B)$) le
système multiplicatif de $\mathcal D$ (resp. $\mathcal A$) associé à
$\mathcal B$ (2.1.7).

\VerdierSemanticAnchor{ch2-2-3-1-a-intersection}
\par\noindent a) On a
$S_{\mathcal D}(\mathcal B)\cap\operatorname{Fl}(\mathcal A)
=S_{\mathcal A}(\mathcal B)$.

\VerdierSemanticAnchor{ch2-2-3-1-b-fully-faithful}
\par\noindent b) Le foncteur canonique
$\mathcal A/\mathcal B\longrightarrow\mathcal D/\mathcal B$ est pleinement
fidèle et injectif sur les objets. L'image de ce foncteur est
$Q(\mathcal A)$, où
$Q\colon\mathcal D\longrightarrow\mathcal D/\mathcal B$ est le foncteur
canonique.

\VerdierSemanticAnchor{ch2-2-3-1-c-quotient-transitivity}
\par\noindent c) Le foncteur canonique
$\mathcal D/\mathcal A\longrightarrow
(\mathcal D/\mathcal B)/(\mathcal A/\mathcal B)$ est un isomorphisme de
catégories.

\VerdierSemanticAnchor{ch2-2-3-1-d-saturation}
\par\noindent d) Si la catégorie $\mathcal A$ est saturée (2.1.6), la
sous-catégorie $Q(\mathcal A)$ est saturée.

\VerdierSemanticAnchor{ch2-2-3-1-inverse-image-setup}
Pour toute sous-catégorie triangulée pleine $\mathcal C$ de
$\mathcal D/\mathcal B$, désignons par $Q^{-1}(\mathcal C)$ la
sous-catégorie pleine de $\mathcal D$ définie par les objets $X$ de
$\mathcal D$ tels que $Q(X)$ soit isomorphe à un objet de la catégorie
$\mathcal C$.

\VerdierSemanticAnchor{ch2-2-3-1-a-bis}
\par\noindent a)$^{\mathrm{bis}}$ La catégorie $Q^{-1}(\mathcal C)$ est une
sous-catégorie triangulée strictement pleine de $\mathcal D$.

\VerdierSemanticAnchor{ch2-2-3-1-b-bis}
\par\noindent b)$^{\mathrm{bis}}$ Le système multiplicatif
$S_{\mathcal D}(Q^{-1}(\mathcal C))$ est composé des morphismes de
$\mathcal D$ dont l'image par $Q$ est un morphisme de
$S_{\mathcal D/\mathcal B}(\mathcal C)$.

\VerdierSemanticAnchor{ch2-2-3-1-c-bis}
\par\noindent c)$^{\mathrm{bis}}$ La sous-catégorie $Q^{-1}(\mathcal C)$ est
saturée si et seulement si la sous-catégorie $\mathcal C$ est saturée.

\VerdierSemanticAnchor{ch2-2-3-1-d-bis}
\par\noindent d)$^{\mathrm{bis}}$ L'application
$\mathcal A\longmapsto Q(\mathcal A)$ est un isomorphisme d'ensembles
ordonnés de l'ensemble des sous-catégories triangulées strictement pleines
de $\mathcal D$ contenant $\widehat{\mathcal B}$ (le noyau du foncteur
$Q$), ordonné par inclusion, dans l'ensemble des sous-catégories
triangulées strictement pleines de $\mathcal D/\mathcal B$, ordonné par
inclusion. L'application inverse est l'application
$\mathcal C\longmapsto Q^{-1}(\mathcal C)$.
\endgroup

\VerdierSemanticAnchor{ch2-2-3-1-proof-ab}
La première assertion est évidente. La deuxième résulte immédiatement de
la description des morphismes dans les catégories
$\mathcal A/\mathcal B$ et $\mathcal D/\mathcal B$ (2.2.1).

\VerdierSemanticAnchor{ch2-2-3-1-proof-c}
Pour démontrer c), il suffit de se reporter aux propriétés universelles des
catégories quotients (2.2.11).

\VerdierSemanticAnchor{ch2-2-3-1-proof-d}
Démontrons d). Quitte à remplacer $\mathcal A$ par sa clôture strictement
pleine, ce qui modifie $Q(\mathcal A)$ par une équivalence, on peut supposer
que $\mathcal A$ est strictement pleine. Le noyau du foncteur
$\mathcal D\longrightarrow\mathcal D/\mathcal A$ est alors la catégorie
$\mathcal A$ (2.2.11). Ce foncteur se factorise en
$\mathcal D\longrightarrow\mathcal D/\mathcal B
\longrightarrow\mathcal D/\mathcal A$. On en déduit que le noyau du
foncteur canonique
$\mathcal D/\mathcal B\longrightarrow\mathcal D/\mathcal A$ est
$Q(\mathcal A)$. Par suite, $Q(\mathcal A)$ est saturée (2.1.16).

\VerdierSemanticAnchor{ch2-2-3-1-proof-a-bis-b-bis}
L'assertion a)$^{\mathrm{bis}}$ est évidente, de même que l'assertion
b)$^{\mathrm{bis}}$.

\VerdierSemanticAnchor{ch2-2-3-1-proof-c-bis-forward}
Supposons que $\mathcal C$ soit saturée. On peut supposer de plus que
$\mathcal C$ est strictement pleine, car si l'on désigne par
$\widehat{\mathcal C}$ la clôture strictement pleine de $\mathcal C$, la
catégorie $\widehat{\mathcal C}$ est strictement pleine saturée et
$Q^{-1}(\widehat{\mathcal C})=Q^{-1}(\mathcal C)$. La catégorie
$\mathcal C$ est alors noyau d'un foncteur exact (2.2.11). Par suite,
$Q^{-1}(\mathcal C)$ est noyau d'un foncteur exact, donc saturée (2.1.16).

\VerdierSemanticAnchor{ch2-2-3-1-p144-boundary}
Supposons que $Q^{-1}(\mathcal C)$ soit
\endgroup
